\documentclass[11pt]{article}
\usepackage[margin=1in]{geometry}
\usepackage[T1]{fontenc}
\usepackage{longtable}
\usepackage{array}
\begin{document}
\section*{FLCR-05 - Typed Boundary Repair}

**Series:** Formalizing LCR, Unifying Standard Models
**Artifact:** formal paper source
**Status:** first-pass rich rewrite; requires final citation and build pass.
**Claim posture:** maximal local claim posture with explicit lane boundaries.

\subsection*{Abstract}

This paper turns failed joins into typed, replayable repair rows. Boundary repair is not a proof of the desired downstream claim; it is a typed exception and proof-obligation discipline that preserves state, coordinate, cause, route, source, target, and falsifier. The main result is an idempotent normalization of failure into legal next-step data.

\subsection*{Keywords}

boundary repair; typed exception; proof obligation; idempotence; schema.

\subsection*{External Reader Orientation}

An outside reviewer should read this paper as a formal-system construction paper. Its local object is **Typed Boundary Repair**. The paper's immediate contribution is: Turns failed joins and boundary mismatches into typed repair data.

The nearest external anchors are finite-state reasoning, typed interfaces, proof-carrying records, conservative extension, and ordinary mathematical citation discipline. LCR terms are therefore internal labels for the construction unless a row explicitly imports an external theorem, citation, dataset, or calibration target.

It does not ask the reader to accept any Standard Model or literal-physics identification at this layer. Such mappings are deferred to the translation papers and must carry their own evidence.

\subsection*{Reviewer Compression}

**Core object.** Typed Boundary Repair.

**Primary result.** This paper contributes the local formal object and its safe downstream-use rule.

**Primary non-result.** This paper does not claim direct identity with Standard Model or physical objects.

**Review strategy.** Evaluate the formal claim rows first, then inspect the receipt/citation/calibration rows that support each claim. Appendices provide routing and implementation context; they should not be read as stronger than their lane permits.

\subsection*{Peer Reviewer Assessment}

Reviewer assumption: the reviewer has been briefed on the FLCR structure, claim lanes, CAM/crystal routing, forge/action surfaces, and the distinction between internal formal results and external physical calibration.

**Recommendation.** major revision, technically promising.

**Review score vector.** orientation=2, claim\_granularity=2, evidence\_visibility=2, falsifier\_visibility=2, journal\_apparatus=1, base\_layer\_boundary=2.

**Formal claim rows reviewed.** 3.

\subsubsection*{Reviewer Strength Finding}

The manuscript is now readable as a bounded formal-system paper: it identifies its local object, separates internal construction from physical interpretation, and exposes claim lanes for review.

\subsubsection*{Major Revision Requests}

\noindent{}-- Convert the strongest proof sketch into numbered definitions, lemmas, and propositions with explicit dependencies.\\
\noindent{}-- Replace placeholder source classes with final citation keys, receipt hashes, or dataset identifiers wherever possible.\\
\noindent{}-- Move repeated pass-derived material into a compact evidence appendix and keep the main body focused on the paper-local argument.\\
\noindent{}-- Add at least one worked example showing the local object, legal operation, receipt, residue, and downstream import rule.\\

\subsubsection*{Minor Revision Requests}

\noindent{}-- Normalize theorem capitalization and punctuation.\\
\noindent{}-- Add final figure/table captions where the text currently describes diagrams schematically.\\
\noindent{}-- Ensure every acronym and local term appears in the paper-local glossary.\\

\subsubsection*{Acceptance Blockers}

\noindent{}-- A reviewer can accept the formal contribution without accepting later physics claims, but only if final citations and receipt identifiers are attached.\\
\noindent{}-- The proof sketch must be promoted into a formal proof or explicitly labeled as a construction argument.\\


\subsection*{1. Contribution And Scope}

\noindent{}-- Defines repair rows as typed constraints rather than proof promotions.\\
\noindent{}-- Installs rejection of untyped failures as a formal schema rule.\\
\noindent{}-- Connects repair rows to graph edges for later causal ledgers.\\
\noindent{}-- Routes physical curvature, dust, oloid, and material readings downstream.\\

This paper belongs to the LCR-native construction band. Physics and Standard Model language, when it appears, is only a deferred mapping cue, analogy, or explicitly bounded calibration target. The paper's base object must remain stable enough for later papers to cite without importing a stronger physical claim by implication.

\subsection*{2. Source Routing}

Primary routed sources:

\noindent{}-- `04\_Boundary\_Repair.md`\\
\noindent{}-- `supplements/04\_INTERNAL\_REASONING\_SUPPLEMENT.md`\\

Cross-corpus evidence classes:

\noindent{}-- `CAM\_CLAIM\_100\_SCORE\_AUDIT.md`\\
\noindent{}-- `NSIT\_TOOL\_CLOSURE\_MATRIX.md`\\
\noindent{}-- `NSIT\_REASONED\_CLOSURE\_PASS.md`\\
\noindent{}-- `PAPER\_REWORKS\_CRYSTAL\_PROJECTION.json`\\
\noindent{}-- standards, glue, window, and node databases where applicable\\
\noindent{}-- notebook-derived review prompts and orbital source manifests\\

\subsection*{3. Definitions}

\noindent{}-- **Repair row.** A structured obligation record with state, coordinate, reason, route, source, target, boundary, and receipt fields.\\
\noindent{}-- **Typed exception.** A failure promoted into a catchable, inspectable, replayable record.\\
\noindent{}-- **Idempotent repair pass.** A normalizer that leaves already-typed repair rows unchanged.\\
\noindent{}-- **Untyped failure.** A failure marker lacking the fields needed for routing or replay.\\

\subsection*{4. Formal Claims}

\noindent\texttt{| Claim | Lane | Statement |}\\
\noindent\texttt{|-------|------|-----------|}\\
\noindent\texttt{| Theorem 5.1 | `receipt\_bound\_internal\_result` | A valid boundary repair converts a failed join into a typed proof-obligation row while preserving the relevant coordinates. |}\\
\noindent\texttt{| Proposition 5.2 | `receipt\_bound\_internal\_result` | Repair is idempotent on already-normalized rows. |}\\
\noindent\texttt{| Protocol 5.3 | `normal\_form\_result` | Untyped failure cannot be consumed by later papers as repair evidence. |}\\

\subsection*{Reviewer Claim Ledger}

This ledger expands the formal-claims table into review actions. It is intended to make proof granularity explicit: what is being claimed, what evidence type can support it, what boundary remains, and what the next review action is.

\noindent\texttt{| Formal row | Type | Claim in review terms | Evidence required | Boundary | Next review action |}\\
\noindent\texttt{|---|---|---|---|---|---|}\\
\noindent\texttt{| Theorem 5.1 | finite/executable result | A valid boundary repair converts a failed join into a typed proof-obligation row while preserving the relevant coordinates | receipt, validator, executable pass, generated artif}\\
\noindent\texttt{| Proposition 5.2 | finite/executable result | Repair is idempotent on already-normalized rows | receipt, validator, executable pass, generated artifact, or finite enumeration | the stated finite/executable domain only |}\\
\noindent\texttt{| Protocol 5.3 | formal construction | Untyped failure cannot be consumed by later papers as repair evidence | definitions, normal-form construction, conversion rule, and downstream-use boundary | internal formal coheren}\\

\subsection*{Claim Granularity And Review Table}

The table below separates the claim types so the paper is reviewable without accepting the whole FLCR vocabulary at once.

\noindent\texttt{| Claim type | What this paper may claim | Acceptance test | What is not claimed by that row |}\\
\noindent\texttt{|---|---|---|---|}\\
\noindent\texttt{| Formal-system result | `FLCR-05` defines or uses **Typed Boundary Repair** as a typed FLCR object with declared inputs, operations, outputs, and residue. | Definitions, formal claims, construction steps, and downstream}\\
\noindent\texttt{| Computational or receipt-bound result | Enumerations, TarPit runs, CAM/crystal routes, verifier rows, and generated artifacts are claims about the stated finite or executable domain. | A named receipt, validator, manif}\\
\noindent\texttt{| Imported theorem or external definition | Imported mathematics remains an external theorem/citation target; this paper may use it only through declared mappings and conservative-extension discipline. | The source theor}\\
\noindent\texttt{| Interpretive bridge or analogy | Analog, workbook, visual, or translation language may explain why the construction is useful. | The analogy preserves the relevant structure and does not promote itself into a theorem. }\\
\noindent\texttt{| Physical or calibration-facing claim | Physical or Standard Model identity is not claimed here unless the paper contains an explicit calibration row; the default status is deferred mapping. | A dataset, citation, calib}\\
\noindent\texttt{| Open obligation or falsifier | A missing derivation, adapter, receipt, dataset, or failed comparator is a named research channel. | The next binding action and the reason the stronger claim is not closed are stated. | }\\

\subsection*{5. Paper-Specific Development}

1. Take an untyped failure and reject it as non-consumable evidence.
2. Normalize the failure into a repair row with state, coordinate, reason, route, source, target, and boundary.
3. Run the normalizer again and confirm that an already typed row is unchanged.
4. Export the row as a proof obligation to the causal ledger, not as the proof of the desired claim.

\subsubsection*{Proof Sketch}

The repair operation is a normalization function on failure records. Records lacking required fields are outside the valid repair language. Records containing the required fields are stable under the operation, so the pass is idempotent. The normalized row is useful precisely because it is replayable and bounded, not because it proves the target.

\subsection*{6. Construction}

The construction is intentionally staged. First, the paper names the finite or
typed object that can be inspected directly. Second, it declares the admissible
operations over that object. Third, it records the receipt or source class that
allows another paper to consume the result. Fourth, it names the residue that
must not be silently promoted.

For this paper, the operative object is `Typed Boundary Repair`. Its safe consumption
rule is:

\begin{verbatim}
source object -> declared operation -> receipt/lane -> downstream import
\end{verbatim}

The unsafe consumption rule is:

\begin{verbatim}
source phrase -> downstream identity claim
\end{verbatim}

That second pattern is explicitly rejected by FLCR-01 and must be rewritten as
a claim-lane promotion with evidence.

\subsection*{7. Evidence And Receipts}

\noindent\texttt{| Evidence source | What it contributes |}\\
\noindent\texttt{|-----------------|---------------------|}\\
\noindent\texttt{| Paper 04 legacy body | Typed boundary repair theorem, verifier, falsifier, repair contract variants. |}\\
\noindent\texttt{| Internal supplement 04 | Typed exception, proof obligation, Hoare-style pre/postcondition framing. |}\\
\noindent\texttt{| CAM Claim 100 Score Audit | Paper 04 scored 93 internally closed. |}\\
\noindent\texttt{| NSIT standards lanes | Claim envelope and content-addressed receipt standards support repair-row normalization. |}\\

\subsection*{Appendix D. Resolved-State Closure Pass}

This pass removes false restrictions on the paper's claim posture. A row is no longer called open merely because a stronger future claim exists. The satisfied lane is closed now; only the stronger claim that lacks its required adapter, comparator, calibration, or falsifier remains as residue.

\subsubsection*{Closed Now}

\noindent\texttt{| Row | Lane | Resolved state | Remaining boundary |}\\
\noindent\texttt{|---|---|---|---|}\\
\noindent\texttt{| paper lift enumeration | `receipt\_bound\_internal\_result` | 8/8 LCR states succeeded; error walls=0; avg TarPit mass=0.514103. | This closes the paper-local finite lift readout, not every stronger physical interpretatio}\\
\noindent\texttt{| boundary repair | `normal\_form\_result` | The typed boundary repair surface is closed as an admission table for legal next-state repair. | The family closure is a structural lane; measured physics still requires destina}\\
\noindent\texttt{| decade-1 tower | `receipt\_bound\_internal\_result` | The decade tower is resolved with avg TarPit mass=0.510370 and mass sum=40.829567. | The decade total is an internal tower metric, not a standalone universal physical }\\
\noindent\texttt{| family-05 cross-lift comparison | `cam\_crystal\_reapplication\_result` | Family tower binds FLCR-05, FLCR-15, FLCR-25, FLCR-35; avg TarPit mass=0.514757; error walls=0. | Cross-lift agreement strengthens routing and comp}\\
\noindent\texttt{| typed repair admission | `normal\_form\_result` | Boundary repair is closed as a typed admission protocol. | Specific physical repair events require process-specific data. |}\\

\subsubsection*{Claim Posture After This Pass}

`FLCR-05` should now be read as a resolved-state contribution for `Typed Boundary Repair` at its declared lane. The paper may state the strongest claim supported by these rows directly. It should reserve caution only for a specifically named stronger claim whose required evidence is absent.

\subsection*{8. Dependencies And Downstream Use}

This paper may be imported by later FLCR papers only through the claim lanes
above. The default downstream consumers are:

\noindent{}-- `FLCR-11` through `FLCR-18` for window, receipt, admission, CA, algebraic,\\
  curvature, residue, and exceptional machinery.
\noindent{}-- `FLCR-28` through `FLCR-30` for CAM/crystal, corpus ribbon, and universal\\
  normal-form intake.
\noindent{}-- `FLCR-31` through `FLCR-40` only after the LCR-native result has been cited\\
  and the translation claim has its own evidence lane.

\subsection*{9. Limitations And Falsifiers}

\noindent{}-- A repair row is an input to future verification, not a proof of the repaired claim.\\
\noindent{}-- Domain-specific repair rows may extend the schema but cannot drop the replay fields.\\
\noindent{}-- Curvature or material interpretations require external equations, data, or calibration receipts.\\

A falsifier for this paper is any example that satisfies the declared input
conditions while violating the stated construction, receipt, or lane boundary.
An interpretation that merely wants a stronger downstream conclusion is not a
falsifier; it is an obligation routed to a later paper.

\subsection*{10. Publication Readiness Tasks}

\noindent{}-- Validator identities and receipt hashes are recorded in the manifest or receipt appendix.\\
\noindent{}-- External theorems require final citation entries before journal submission.\\
\noindent{}-- Every theorem, proposition, and protocol must retain a claim lane and evidence boundary.\\
\noindent{}-- The Markdown, LaTeX, and PDF forms are generated from the same canonical source.\\

\subsection*{Dual Positioning: Story And Formal Carrier}

This paper must be read in two synchronized ways. Sequentially, it explains why
the next state exists in the corpus story. Formally, it binds that state to
accepted mathematics, IRL formalism, code receipts, validators, CAM/crystal
lookups, and claim-lane boundaries.

Assigned original ribbon role(s): `04`/boundary\_action.

\noindent\texttt{| Original slot | Family | Lift | Role | Proof form |}\\
\noindent\texttt{|---|---|---:|---|---|}\\
\noindent\texttt{| `04` | boundary\_action | 0 | order-1 slot-4: type the boundary, repair row, or curvature-demand condition | typed boundary row and next-state admissibility |}\\

The formal obligation is to state the strongest valid claim available for this
paper without letting interpretation silently change the addressed object. Any
author interpretation belongs beside the formalism as a declared relabeling,
bridge, or analog, and must be accompanied by the evidence lane that makes the
claim consumable by later papers.

\subsection*{State-Bound Proof Contract}

This paper receives: state emitted by prior slot 03 and reopened at original slot 04 (Boundary Repair).

It must produce: formal result of original slot 04 plus the same-family lift path toward slot 14.

Semantic transition: type the boundary condition that decides whether the state repairs, reflects, curves, or routes onward.

Accepted formalism targets to bind in the journal rewrite:

\noindent{}-- boundary value problems\\
\noindent{}-- typed interfaces and admissibility conditions\\
\noindent{}-- topological boundary/interior distinction\\
\noindent{}-- falsifier rows for failed promotions\\

\noindent\texttt{| Slot | Family | Lift | Produced proof form |}\\
\noindent\texttt{|---|---|---:|---|}\\
\noindent\texttt{| `04` | boundary\_action | 0 | typed boundary row and next-state admissibility |}\\

The formal carrier uses these accepted tools while preserving interpretive
claims as explicit relabelings, correspondences, or bridges. Claim promotion is admissible only when a receipt, standard citation,
CAM/crystal reapplication, normal-form result, calibration datum, or falsifier
boundary is attached.

\subsection*{Provenance And Crystal Evidence Integration}

This section integrates prior work, CAM/crystal projections, NSIT tooling,
standards-window evidence, and routed supplements as provenance-bearing
evidence. Source material is not treated as manuscript text; it is bound into
the paper only through claim lanes, receipts, validators, normal forms,
calibration rows, or explicit falsifier boundaries.

\subsubsection*{Expansion Rule For This Slot}

Use past work to type the boundary: admitted, repaired, reflected, calibration-required, falsified, or routed onward.

\subsubsection*{Primary Evidence Inputs}

\noindent\texttt{| Source | Placement reason |}\\
\noindent\texttt{|---|---|}\\
\noindent\texttt{| `04\_Boundary\_Repair.md` | primary assigned source for the paper's formal spine |}\\
\noindent\texttt{| `supplements/04\_INTERNAL\_REASONING\_SUPPLEMENT.md` | paper-local reasoning supplement contributing to definitions, proof sketch, and limitations |}\\

\subsubsection*{Secondary And Orbital Evidence Inputs}

\noindent\texttt{| Source | Placement reason |}\\
\noindent\texttt{|---|---|}\\
\noindent\texttt{| `supplements/RECEIPT\_VERIFIER\_CATALOG.md` | cross-cutting supplement; paper-relevant rows, receipts, and guard language are bound through evidence lanes |}\\
\noindent\texttt{| `supplements/INTERNAL\_REASONING\_PYRAMID\_INDEX.md` | cross-cutting supplement; paper-relevant rows, receipts, and guard language are bound through evidence lanes |}\\
\noindent\texttt{| `supplements/INTERNAL\_REASONING\_5P\_WINDOW\_SUPPLEMENT.md` | cross-cutting supplement; paper-relevant rows, receipts, and guard language are bound through evidence lanes |}\\
\noindent\texttt{| `supplements/INTERNAL\_REASONING\_7P\_WINDOW\_SUPPLEMENT.md` | cross-cutting supplement; paper-relevant rows, receipts, and guard language are bound through evidence lanes |}\\
\noindent\texttt{| `supplements/INTERNAL\_REASONING\_9P\_WINDOW\_SUPPLEMENT.md` | cross-cutting supplement; paper-relevant rows, receipts, and guard language are bound through evidence lanes |}\\
\noindent\texttt{| `CAM\_CLAIM\_100\_SCORE\_AUDIT.md` | series-wide audit/score/closure evidence; import paper-specific rows and leave global context outside the body |}\\
\noindent\texttt{| `NSIT\_TOOL\_CLOSURE\_MATRIX.md` | series-wide audit/score/closure evidence; import paper-specific rows and leave global context outside the body |}\\
\noindent\texttt{| `NSIT\_REASONED\_CLOSURE\_PASS.md` | series-wide audit/score/closure evidence; import paper-specific rows and leave global context outside the body |}\\
\noindent\texttt{| `ORBITAL\_DATA\_CROSS\_POPULATION\_AUDIT.md` | series-wide audit/score/closure evidence; import paper-specific rows and leave global context outside the body |}\\
\noindent\texttt{| `AGENT\_CRYSTAL\_WORK\_HARVEST.md` | series-wide audit/score/closure evidence; import paper-specific rows and leave global context outside the body |}\\
\noindent\texttt{| `PAPER\_REWORKS\_CRYSTAL\_PROJECTION.json` | series-wide audit/score/closure evidence; import paper-specific rows and leave global context outside the body |}\\
\noindent\texttt{| `NEW\_PAPER\_NEEDS.md` | route relevant obligations into this paper's burden table: NP-05 Receipt Ecology And 32-Paper Causal Graph |}\\
\noindent\texttt{| Additional source routes | Complete routing is maintained in the source-placement appendix |}\\

\subsubsection*{Crystal And Standards Evidence To Bind}

\noindent{}-- Paper Reworks crystal projection: 33 paper markdown files, 9 supplements, 5 queues, 6 lattice-forge unification artifacts, 140 faces, 140 vignettes, and 280 CAM nodes.\\
\noindent{}-- Standards-window intake: 192/192 corpus conformance verdicts PASS; 140/142 obligations have candidate routes; 2/4/8/16/32 window lattice is available for cross-paper reads.\\
\noindent{}-- Agent/crystal harvest: agent-generated memory, runtime standards starter, current runtime build, repo-harvest CAM, notebook-derived bundles, and downloaded crystal payloads are orbital evidence surfaces.\\
\noindent{}-- NSIT inventory baseline: 220 validators, 1786 receipt/data artifacts, 1137 formal-paper-like artifacts, 114 supplements, and 86 memory/CAM artifacts.\\

\subsubsection*{Paper-Specific CAM/Score Rows}

\noindent\texttt{| 04 | 93 | internally closed | typed repair constraint, falsifier, Paper 05 payload intake | shared ledger population and domain examples |}\\

\subsubsection*{Paper-Specific NSIT Closure Rows}

No direct NSIT row matched the assigned legacy papers in the first-pass matrix; validator matching by object name remains a receipt-attachment requirement.

\subsubsection*{Source Integration Summary}

The legacy source and internal reasoning supplement are treated as source
evidence rather than manuscript prose. Their contributions enter this paper as
claim-lane entries, receipt or validator references, finite-domain statements,
and limitation boundaries. Speculative or cross-domain interpretations remain
separate from closed local results until an explicit adapter, citation,
normal-form derivation, calibration row, or falsifier is attached.
\subsection*{Claim Binding Queue}

This queue records the evidence discipline for claim strengthening. It separates claims that already have support from residues that remain in falsifier or open-obligation lanes.

\noindent\texttt{| Queue | Promote now | Bind next | Residue |}\\
\noindent\texttt{|---|---|---|---|}\\
\noindent\texttt{| none directly assigned | Source routing and claim lanes identify paper-local binding actions. | Verifier catalogs and CAM/crystal routes provide object-name evidence candidates. | Unbound claims remain in falsifier/ope}\\

\subsubsection*{Binding Discipline}

Each queue item is represented as a delta:

\begin{verbatim}
local claim at paper time
-> attached later source/tool/receipt/theorem
-> revised representation
-> invariant boundary and remaining residue
\end{verbatim}

This preserves the sequential story while allowing later tools, crystals, and
standard formalisms to return through the proper lane.

\subsection*{Claim Delta Ledger}

This ledger applies the paper's claim-binding queues as explicit back-application
deltas. It keeps the sequential paper story intact while allowing later
receipts, standard formalisms, CAM/crystal routes, and validators to sharpen the
claim.

\noindent\texttt{| Delta | Local claim at paper time | Later evidence | Revised representation | Boundary kept |}\\
\noindent\texttt{|---|---|---|---|---|}\\
\noindent\texttt{| Search By Object Name | No direct reasoned-closure queue was assigned from the first queue map. | Search source routing, verifier catalog, CAM/crystal routes, and paper-local object names. | Create a specific binding q}\\

The revised representation records the strongest claim supported by the current evidence package. Promotion into theorem or proposition language occurs only when the named evidence is attached in the manifest or workbook replay.

\subsection*{Source-Backed Evidence Summary}

Routed archive and mirror cards are treated as provenance summaries rather than
body text. They identify candidate receipts, validators, theorem registries, or
supplemental source routes. A card strengthens this paper only when its
contribution is bound to a claim lane, evidence object, and boundary.

\noindent\texttt{| Evidence card | Score | Publication contribution | Boundary |}\\
\noindent\texttt{|---|---:|---|---|}\\
\noindent\texttt{| CQE-paper-04.25: Paper 4.25 - Toolkit for Boundary Repair | 81 | Paper 4.25 describes the review tools for boundary repair. The tools help a reader construct and inspect repair rows. They do not add claims beyond the P}\\
\noindent\texttt{| FINAL\_FORMAL\_PAPER: Complete Formal Claims: The Folded Strand | 75 | We present the complete closed-form claim set of the CQE\_CMPLX corpus: - 33 papers = 33 folding operations on a self-complementary strand - 144 tools}\\
\noindent\texttt{| SUMMARY-V-32-Theorems-Registry: Summary Paper V — The 32 Theorems in One Table: Closed-Form Registry | 75 | This paper is the **complete closed-form registry of all 32 theorems** in the CQE\_CMPLX corpus. Each theorem i}\\
\noindent\texttt{| CQE-paper-02.50: Paper 2.50 - Correction Surface Claim Contract | 73 | Paper 2.50 lets later papers import the correction surface honestly. It keeps the finite theorem available while preserving the distinction between}\\
\noindent\texttt{| FORMAL: Paper 01 — Side-Flip SU(2) Doublet (T\_BIJECTIVE) | 73 | **PROVEN by structural identification** on the T3 chart/J₃(O) isomorphism. | Source-backed support; promotion requires the paper-local claim lane and evid}\\
\noindent\texttt{| Additional evidence cards | 23 total | The complete card inventory is maintained in the archive evidence matrix. | Cards are source-discovery aids until bound to paper-local evidence. |}\\

\subsection*{Journal Apparatus}

\subsubsection*{Article Type And Intended Venue Fit}

This manuscript is written as a formal-methods and mathematical-systems paper
inside the series *Formalizing LCR, Unifying Standard Models*. Its intended
reader is a technical reviewer who needs claim boundaries, reproducibility
routes, and cross-paper dependencies to be visible without relying on private
context.

\subsubsection*{Structured Contribution Statement}

\noindent\texttt{| Item | Statement |}\\
\noindent\texttt{|---|---|}\\
\noindent\texttt{| Publication ID | `FLCR-05` |}\\
\noindent\texttt{| Legacy source slot(s) | `04` |}\\
\noindent\texttt{| Ribbon role | order-1 slot-4: type the boundary, repair row, or curvature-demand condition |}\\
\noindent\texttt{| Proof form | typed boundary row and next-state admissibility |}\\
\noindent\texttt{| Received state | state emitted by prior slot 03 and reopened at original slot 04 (Boundary Repair) |}\\
\noindent\texttt{| Produced state | formal result of original slot 04 plus the same-family lift path toward slot 14 |}\\

\subsubsection*{Claim-Evidence Table}

\noindent\texttt{| Claim | Lane | Evidence to attach | Boundary |}\\
\noindent\texttt{|---|---|---|---|}\\
\noindent\texttt{| Theorem 5.1 | `receipt\_bound\_internal\_result` | Receipt, source card, validator, citation, or CAM route named in the paper manifest | A valid boundary repair converts a failed join into a typed proof-obligation row whi}\\
\noindent\texttt{| Proposition 5.2 | `receipt\_bound\_internal\_result` | Receipt, source card, validator, citation, or CAM route named in the paper manifest | Repair is idempotent on already-normalized rows. |}\\
\noindent\texttt{| Protocol 5.3 | `normal\_form\_result` | Receipt, source card, validator, citation, or CAM route named in the paper manifest | Untyped failure cannot be consumed by later papers as repair evidence. |}\\

\subsubsection*{Figure Plan}

\noindent\texttt{| Figure | Purpose | Caption |}\\
\noindent\texttt{|---|---|---|}\\
\noindent\texttt{| FLCR-05-F1 | Typed boundary/admission diagram | Schematic showing how `FLCR-05` turns its received state into the produced state while preserving claim lanes and residue boundaries. |}\\
\noindent\texttt{| FLCR-05-F2 | Evidence routing map | Diagram of source papers, archive cards, CAM/crystal routes, validators, and workbook replay paths used by this manuscript. |}\\
\noindent\texttt{| FLCR-05-F3 | Same-family lift context | Diagram placing this paper in the original `00-79` ribbon family and showing predecessor/successor slots. |}\\

\subsubsection*{In-Text Figure FLCR-05-F1: Typed boundary/admission diagram}

\begin{verbatim}
received state
  -> Typed boundary/admission diagram
  -> formal claim lane
  -> evidence object / receipt / citation
  -> produced state
  -> remaining residue
\end{verbatim}

\subsubsection*{Table Plan}

\noindent\texttt{| Table | Purpose |}\\
\noindent\texttt{|---|---|}\\
\noindent\texttt{| FLCR-05-T1 | Boundary verdict table |}\\
\noindent\texttt{| FLCR-05-T2 | Claim-lane/evidence-boundary table |}\\
\noindent\texttt{| FLCR-05-T3 | Archive evidence card and source-backed upgrade table |}\\
\noindent\texttt{| FLCR-05-T4 | Workbook replay and falsifier table |}\\

\subsubsection*{Worked Example}

Classify a candidate as admitted, repaired, boundary, rejected, or calibration-required. The example must name the input object, the operation,
the evidence object, the allowed revised claim, and the remaining boundary.

\subsubsection*{Nomenclature And Glossary}

\noindent\texttt{| Term | Paper-local meaning |}\\
\noindent\texttt{|---|---|}\\
\noindent\texttt{| CAM | Defined by this paper as part of the `boundary\_action` proof role; final definition is recorded in the paper-local glossary. |}\\
\noindent\texttt{| admission | Defined by this paper as part of the `boundary\_action` proof role; final definition is recorded in the paper-local glossary. |}\\
\noindent\texttt{| boundary | Defined by this paper as part of the `boundary\_action` proof role; final definition is recorded in the paper-local glossary. |}\\
\noindent\texttt{| claim lane | Defined by this paper as part of the `boundary\_action` proof role; final definition is recorded in the paper-local glossary. |}\\
\noindent\texttt{| crystal | Defined by this paper as part of the `boundary\_action` proof role; final definition is recorded in the paper-local glossary. |}\\
\noindent\texttt{| falsifier | Defined by this paper as part of the `boundary\_action` proof role; final definition is recorded in the paper-local glossary. |}\\
\noindent\texttt{| receipt | Defined by this paper as part of the `boundary\_action` proof role; final definition is recorded in the paper-local glossary. |}\\
\noindent\texttt{| repair row | Defined by this paper as part of the `boundary\_action` proof role; final definition is recorded in the paper-local glossary. |}\\

\subsubsection*{Appendix FLCR-05-A: Evidence Package}

The appendix records:

1. Legacy source excerpts or source-card references used in the final claim ledger.
2. Receipt and verifier paths, including pass/fail/partial status.
3. CAM/crystal query or projection rows used by the paper, with source identity and boundary.
4. Standards-window and NSIT evidence used for conformance or routing.
5. Falsifier, calibration, or external-data rows that remain unresolved.

\subsubsection*{Data And Code Availability}

The publication package contains the canonical Markdown sources, manifests, source-routing matrices, validation reports, and generated PDF/TeX outputs.

Code and verifier references are cited through manifests, archive evidence cards, receipt catalogs, or workbook replay tables. External datasets or
standards citations must be attached before any calibration-bound claim is
promoted.

\subsubsection*{Reviewer Checklist}

\noindent\texttt{| Check | Reviewer question |}\\
\noindent\texttt{|---|---|}\\
\noindent\texttt{| Claim lane | Does every strong claim declare its lane and evidence type? |}\\
\noindent\texttt{| Boundary | Does the paper preserve what it does not prove? |}\\
\noindent\texttt{| Reproducibility | Is there a receipt, verifier, source card, citation, or replay table for the claim? |}\\
\noindent\texttt{| Cross-paper import | Does the paper cite the prior FLCR/OPR slot before consuming its result? |}\\
\noindent\texttt{| Interpretation | Does the analog layer preserve the addressed state while changing labels/access paths? |}\\

\subsubsection*{Declarations}

No human-subjects, clinical, or animal-research protocol is asserted by this
paper. External empirical claims require separate dataset, calibration, or
domain-validation attachments. Competing-interest and funding statements are recorded in the submission metadata.

\subsection*{11. Publication Integration Addendum}

\subsubsection*{11.1 Role In The Series}

`FLCR-05` (`Typed Boundary Repair`) occupies the `boundary\_action` role at lift depth `0`.
The paper receives state emitted by prior slot 03 and reopened at original slot 04 (Boundary Repair). It produces formal result of original slot 04 plus the same-family lift path toward slot 14. The state transition is:
type the boundary condition that decides whether the state repairs, reflects, curves, or routes onward.

The paper-local proof form is:

\begin{verbatim}
typed boundary row and next-state admissibility
\end{verbatim}

The integration result is a typed boundary/admission table for legal next-state repair. This addendum records how that result is
consumed by the publication series without relying on editorial instructions or
private working context.

\subsubsection*{11.2 Formal Carrier}

\noindent\texttt{| Formal carrier | Function |}\\
\noindent\texttt{|---|---|}\\
\noindent\texttt{| boundary value problems | Formal carrier for the paper-local result. |}\\
\noindent\texttt{| typed interfaces and admissibility conditions | Formal carrier for the paper-local result. |}\\
\noindent\texttt{| topological boundary/interior distinction | Formal carrier for the paper-local result. |}\\
\noindent\texttt{| falsifier rows for failed promotions | Formal carrier for the paper-local result. |}\\

The LCR, CAM, crystal, forge, and analog vocabulary functions as a typed access
layer over these carriers. A relabeling is admissible only when the addressed
object, evidence lane, boundary, and downstream import rule remain attached.

\subsubsection*{11.3 Claim Ledger}

\noindent\texttt{| Claim | Lane | Statement |}\\
\noindent\texttt{|---|---|---|}\\
\noindent\texttt{| Theorem 5.1 | `receipt\_bound\_internal\_result` | A valid boundary repair converts a failed join into a typed proof-obligation row while preserving the relevant coordinates. |}\\
\noindent\texttt{| Proposition 5.2 | `receipt\_bound\_internal\_result` | Repair is idempotent on already-normalized rows. |}\\
\noindent\texttt{| Protocol 5.3 | `normal\_form\_result` | Untyped failure cannot be consumed by later papers as repair evidence. |}\\

These claims are consumed at their stated lane. A stronger downstream statement
creates a new claim envelope with its own evidence object.

\subsubsection*{11.4 Evidence Package}

\noindent\texttt{| Evidence class | Routed items | Status |}\\
\noindent\texttt{|---|---|---|}\\
\noindent\texttt{| Legacy sources | `04\_Boundary\_Repair.md`, `supplements/04\_INTERNAL\_REASONING\_SUPPLEMENT.md` | routed evidence |}\\
\noindent\texttt{| Evidence inputs | `CAM\_CLAIM\_100\_SCORE\_AUDIT.md`, `NSIT\_TOOL\_CLOSURE\_MATRIX.md`, `NSIT\_REASONED\_CLOSURE\_PASS.md`, `PAPER\_REWORKS\_CRYSTAL\_PROJECTION.json` | routed evidence |}\\
\noindent\texttt{| CAM/crystal sources | `PAPER\_REWORKS\_CRYSTAL\_PROJECTION.json`, `AGENT\_CRYSTAL\_WORK\_HARVEST.json`, `runtime standards artifact`, `notebook-derived source bundle` | routed evidence |}\\
\noindent\texttt{| Standards-window inputs | `window\_summary.json`, `node DBs`, `window DBs` | routed evidence |}\\
\noindent\texttt{| Receipts | external requirement | not attached in this source package |}\\
\noindent\texttt{| Validators | external requirement | not attached in this source package |}\\

Standards-window, runtime, and notebook-derived artifacts are treated
as evidence-routing surfaces. They support source discovery, conformance
checking, and reapplication, but they do not replace citations, receipts,
normal-form derivations, calibration rows, or falsifiers.

\subsubsection*{11.5 Results Representation}

The paper's results section is organized around five publication objects:

1. the exact object acted on at this slot;
2. the admissible operation or transformation;
3. the receipt, enumeration, derivation, lookup, or external theorem supporting
   the operation;
4. the residue or boundary left after the result;
5. the downstream import rule.

\subsubsection*{11.6 Figures And Tables}

\noindent\texttt{| Figure | Role |}\\
\noindent\texttt{|---|---|}\\
\noindent\texttt{| FLCR-05-F1: Typed boundary/admission diagram | Visualizes the proof object, routing, or boundary. |}\\
\noindent\texttt{| FLCR-05-F2: Evidence routing map | Visualizes the proof object, routing, or boundary. |}\\
\noindent\texttt{| FLCR-05-F3: Same-family lift context | Visualizes the proof object, routing, or boundary. |}\\

\noindent\texttt{| Table | Role |}\\
\noindent\texttt{|---|---|}\\
\noindent\texttt{| FLCR-05-T1: Boundary verdict table | Records claim lanes, evidence, sources, or falsifiers. |}\\
\noindent\texttt{| FLCR-05-T2: Claim-lane/evidence-boundary table | Records claim lanes, evidence, sources, or falsifiers. |}\\
\noindent\texttt{| FLCR-05-T3: Archive evidence card and source-backed upgrade table | Records claim lanes, evidence, sources, or falsifiers. |}\\
\noindent\texttt{| FLCR-05-T4: Workbook replay and falsifier table | Records claim lanes, evidence, sources, or falsifiers. |}\\

\subsubsection*{11.7 Limits And Falsifiers}

The paper proves only the object and domain declared in its claim ledger.
Physical, Standard Model, observer, calibration, or synthesis claims require
the additional lane evidence stated by their destination papers. CAM/crystal
and forge language is operational only when represented as an address, lookup,
receipt, validator, or reapplication path.
\end{document}
